\section{Heap Page：页结构、校验与记录切片遍历}

对应实现：\fpath{newdb/src/heap_page_waterfall.cpp}

\subsection{关键代码段：追加记录（slot 表 + free space）}

追加一条记录的核心逻辑是“payload 向前增长 + slot 表向后增长”，并用 header 维护两者之间的空洞：

\begin{lstlisting}[language=C++,caption={append\_encoded\_record 的关键步骤（节选）}]
const std::size_t need = rec_len + sizeof(std::uint16_t);
if (cur_free_size < need) {
    return false;
}
const unsigned short rec_offset = cur_free_space;
std::memcpy(page + rec_offset, encoded, rec_len);

const unsigned short new_num = cur_num + 1;
std::uint16_t* new_slot = slot_ptr(page, new_num);
*new_slot = rec_offset;

const unsigned short slot_region_begin =
    static_cast<unsigned short>(reinterpret_cast<unsigned char*>(new_slot) - page);
const unsigned short new_free_space = static_cast<unsigned short>(rec_offset + rec_len);
const unsigned short new_free_size =
    static_cast<unsigned short>(slot_region_begin - new_free_space);

header->set_num_records(new_num);
header->set_free_space(new_free_space);
header->set_free_size(new_free_size);
\end{lstlisting}

细节解读：

\begin{itemize}
  \item \textbf{为何 need 要加 2 字节}：每条记录除了 payload，还要在 slot 表写入一个 \texttt{uint16\_t} 的偏移值。
  \item \textbf{slot\_ptr 的语义}：slot 表位于页尾，\texttt{slot\_ptr(page, n)} 返回“n 条记录时，新 slot 应写入的位置”。
  \item \textbf{free\_size 的计算}：用“slot 表起始地址”减去“payload 已写到的位置”，保持两端不相撞。
\end{itemize}

\subsection{页内布局（逻辑视图）}

Heap Page 使用 Waterfall 的 \texttt{page\_header\_t}，并维护三类区域：

\begin{itemize}
  \item \textbf{Header}：记录页类型、record 数、free space 指针、free size、CRC32 等；
  \item \textbf{Payload 区}：自前向后增长，连续写入每条 record 的 payload（不显式存长度）；
  \item \textbf{Slot 表}：自页尾向前增长，每新增一条记录，写入一个 \texttt{uint16\_t} 的 payload 偏移。
\end{itemize}

新增记录时，需要的空间为：

\[
\text{need} = \text{payload\_len} + \text{sizeof(uint16\_t)}
\]

若 \texttt{free\_size < need} 则追加失败，调用方通常会新开一页再试。

\subsection{CRC32 校验}

页校验由 \texttt{page\_base::compute\_checksum()} 计算，header 保存的 \texttt{crc32} 与之相等则通过。
装载时校验失败的页会被跳过（容错策略），写入时每次落盘前会更新 CRC32。

\subsection{记录遍历：为什么要对 slot 排序}

Slot 表中记录的“slot 下标顺序”与 payload 的“物理偏移升序”并不一致（例如新写入的 slot 在更靠近页尾的位置）。
而 record 的长度并未显式存储，因此遍历时必须按 payload 偏移升序来恢复每条记录长度：

\begin{enumerate}
  \item 读取 header 的 record 数 \texttt{num} 与 free space 指针 \texttt{free\_space}；
  \item 取到 slot 表全部偏移数组 \texttt{all\_slots[0..num-1]}；
  \item 构造 \texttt{order}，按 \texttt{all\_slots[order[k]]} 升序排序；
  \item 第 \(k\) 条 record 的起点为 \(\text{all\_slots[order[k]]}\)，终点为下一条起点（或 \texttt{free\_space}），长度为两者差。
\end{enumerate}

该算法的时间复杂度为 \(O(num \log num)\)（页内排序），但页内 record 数通常受限于页大小，因此在整体扫描中表现稳定。

\subsection{关键代码段：walk\_record\_slices 的“排序 + 差分求长度”}

\begin{lstlisting}[language=C++,caption={walk\_record\_slices 的核心思路（节选）}]
// 1) order = [0..num-1]
std::sort(order.begin(), order.end(), [&](unsigned short a, unsigned short b) {
    return all_slots[a] < all_slots[b];
});

// 2) rec_len = next_offset - rec_offset
const unsigned short rec_offset = all_slots[slot_index];
const unsigned short next_offset =
    (k + 1 < num) ? all_slots[order[k + 1]] : free_space;
const std::size_t rec_len = static_cast<std::size_t>(next_offset - rec_offset);
slice.payload = page + rec_offset;
slice.payload_length = rec_len;
\end{lstlisting}

为什么这个设计很关键：

\begin{itemize}
  \item \textbf{不存长度，靠差分恢复}：页内不为每条记录保存长度，节省空间，但要求遍历时能从“偏移序列”恢复长度。
  \item \textbf{必须排序}：如果直接按 slot 下标遍历，\texttt{next\_offset} 不一定比当前 offset 大，差分会变负或越界。
  \item \textbf{边界使用 free\_space}：最后一条记录的终点就是 header 维护的 free space 指针。
\end{itemize}

